\section{The scale of the arithmetic comparison}\label{sec:comparison}

The analytic criterion leaves us with a precise arithmetic task. Its family of
Dirichlet characters is fixed, whereas each auxiliary abelian extension brings its own
characters into the prime sum. Averaging over a subgroup of
translates projects onto characters that are trivial on that
subgroup. We identify the resulting Fourier coefficients and
separate the modes involving the two auxiliary quadratic
characters. Once the estimate is divided by the total weight,
the slope errors supplied by
Section~\ref{subsec:fixed-weight-absorption} give an upper bound
for \(G_{\mathcal K}\).

\begin{definition}[Translate group and weight scale]\label{def:translate-scale}
Fix \(x\ge2\), \(\epsilon>0\) and the primitive target \(\psi\) of
order \(o\). Here \(q,d\) denote the two auxiliary primes of
Section~\ref{sec:pair-selection}; the conductor of the target is \(f_\psi\).
Let \(H\) be the Galois group of the auxiliary abelian field over \(\mathbb Q\)
and \(\widehat H=\operatorname{Hom}(H,\mathbb C^\times)\) its character group.
Order the two quadratic coordinates by \(d\) first and \(q\) second.
For \(p\in\{d,q\}\) define \(\epsilon_p(h)\in\mathbb Z/2\mathbb Z\) by
\(h(\sqrt{-p})=(-1)^{\epsilon_p(h)}\sqrt{-p}\).
Let \(A\) be the joint image of \(\psi\) and the Galois action on the real
cyclotomic field of \(2\)-power conductor in Definition~\ref{def:family-original-host}.
Write \(\pi_\psi:A\twoheadrightarrow\psi(U_{f_\psi})\) for projection to the value of \(\psi\). Thus
\[
 H=(\mathbb Z/2\mathbb Z)^2\times A,\qquad
 A_0=\ker\pi_\psi,\quad [A:A_0]=o,\qquad
 \widehat H_0=\{\omega\in\widehat H:\omega|_{0\times A_0}=1\}.
\]

In the compositum \(E\) of \eqref{eq:host-character-compositum}, the two
odd primes \(d,q\) are distinct and prime to the conductor of the joint
factor. Ramification at each prime distinguishes its quadratic factor
from the compositum of the others. In these product coordinates put
\[
 \begin{aligned}
 P&=\operatorname{Gal}(E/\mathbb Q(\sqrt{-d}))
       =\{0\}\times\mathbb Z/2\mathbb Z\times A,\\
 \omega_P(h)&=(-1)^{\epsilon_d(h)},&
 e_1(s)&=(1,0,s)\quad(s\in A).
 \end{aligned}
\]
Inertia at \(d\) acts trivially on the other two factors and surjects
onto the ramified quadratic factor. Hence
\[
 I_d=\{0,e_1(0)\},\qquad H\setminus P=e_1(0)+P.
\]

Put \(D=2^e\). The automorphism \(\zeta\mapsto\zeta^5\) identifies
\(\operatorname{Gal}(C_e/\mathbb Q)\) with the additive group
\(\mathbb Z/D\mathbb Z\). Indeed successive squaring gives
\(v_2(5^{2^i}-1)=i+2\), and the powers of 5 represent the units
modulo \(2^{e+2}\) up to sign. Write \(\rho:A\twoheadrightarrow
\mathbb Z/D\mathbb Z\) for restriction. The joint image embeds in
\(\psi(U_{f_\psi})\times\mathbb Z/D\mathbb Z\). On \(A_0\) its first
coordinate is trivial. Thus \(\rho|_{A_0}\) is injective. If
\(t_0=[\mathbb Z/D\mathbb Z:\rho(A_0)]\) and \(n_0=|A_0|\), then
\begin{equation}\label{eq:dyadic-coordinate-orders}
 \rho(A_0)=t_0\mathbb Z/D\mathbb Z,\qquad
 n_0=D/t_0=n/o,\qquad
 \mathbb Z/n_0\mathbb Z\xrightarrow{\ s\mapsto t_0s\ }\rho(A_0).
\end{equation}

The last map is an isomorphism. For \(\psi=1\), one has \(o=t_0=1\)
and \(n=D\). In general \(n\) counts pairs of coordinates.
Let \(T\subseteq A\) index the top weights and \(T_2\subseteq A_0\)
the weights of size approximately \(b>0\) in the other levels.
For \(m\) lifts use the normalization
\[
 n=|A|,\qquad n_T=|T|,\qquad n_2=|T_2|,\qquad
 \delta_2=n_2/n,\qquad Q=n_Tmb.
\]
\end{definition}

\begin{definition}[Centered contribution and omitted characters]\label{def:translate-errors}
For \(1\le j\le m\), \(t\in T\), let \(\Delta_{j,t}:A_0\to\mathbb R\)
be the local prime sum minus its centering at lift \(j\) and top index \(t\),
both computed with the same integral weights. To specify the difference before
Fourier expansion, use the quadratic subgroup \(P\), its character
\(\omega_P\), the top weight \(a=a_t\) and the weights
\(w_g=w_g(j,t,u)\) of Lemma~\ref{hyp:inertia-fourier}.

For each rational prime \(p\) write \(I_p\subset H\) for inertia and
\(F_p\in H\) for any lift of arithmetic Frobenius modulo \(I_p\).
Put \((z)_+=\max(z,0)\) and \(f_P=|P|\). The prime sum is
\[
 V_{j,t}(u)=\sum_{g\notin P}\sum_{p,\nu\ge1}
    \frac{2\log p}{|I_p|}\,\mathbf1_{g\in\nu F_p+I_p}
               (a/p^\nu-w_g)_+ ,
\]
where the inner sum runs over primes \(p\) and integers \(\nu\ge1\).
The positive part makes it finite. With \(\gamma\) as in
Definition~\ref{def:smoothed-centering}, its centering and difference are
\[
 \begin{aligned}
 C_{j,t}(u)&=a\left(\log a-1-\gamma
   -f_P^{-1}\sum_{g\notin P}\log w_g
   +\RePart\frac{L'(1,\chi_{\omega_P})}{L(1,\chi_{\omega_P})}\right),\\
 \Delta_{j,t}(u)&=V_{j,t}(u)-C_{j,t}(u).
 \end{aligned}
\]
Formula~\eqref{eq:color-balance} identifies this finite arithmetic
expression with a sum over characters.

Choose a modulus \(N\) for which the auxiliary field is contained in
\(\mathbb Q(\mu_N)\). Let \(\operatorname{Art}_N:U_N\twoheadrightarrow H\)
be restriction from the cyclotomic Galois group. For a prime \(p\nmid N\),
\(\operatorname{Art}_N([p]_N)\) is arithmetic Frobenius, the restriction
of \(\zeta\mapsto\zeta^p\) on \(N\)-th roots of unity.
For a character \(\omega\in\widehat H\), the primitive Dirichlet character
inducing \(\omega\circ\operatorname{Art}_N\) is
\(\chi_\omega=(\omega\circ\operatorname{Art}_N)^*\), as in Lemma~\ref{hyp:inertia-fourier}.
The absolute contribution from the omitted characters is
\[
 \Sigma(x;\operatorname{Art}_N,A_0,\psi)
   =\sum_{\substack{\omega\in\widehat H_0\\
              \chi_\omega\notin\mathcal K(\psi)}} |G_{\chi_\omega}(x)|.
\]
Write \(\Sigma(x)\) when the Artin quotient, subgroup and target character are fixed.
\end{definition}

For a translate \(u\in A_0\), let \(S_{\rm prime}(u)\) be the sum of the
local prime sums over \(j,t\), and \(S_{\rm cent}(u)\) the sum of their
centering terms. If \(E\) bounds their difference for every translate, then
\[
 \begin{aligned}
 S_{\rm prime}(u)-S_{\rm cent}(u)
   &=\sum_{j=1}^{m}\sum_{t\in T}\Delta_{j,t}(u)\le E
       &&(u\in A_0),\\
 |\widehat H_0|&=|H/(0\times A_0)|=4o .
 \end{aligned}
\]
Take the pair \((q,d)\) supplied by Proposition~\ref{prop:quadratic-prime-pair}
for \([Y,Y^K]\), with \(Y\) above its threshold for the fixed modulus, \(K\) and \(\epsilon\).
For every \(x\) in this interval the proposition bounds each omitted twist by
\(x^{1/2+\epsilon/2}\). In particular,
\[
 \Sigma(x)\le 4o\,x^{1/2+\epsilon}\qquad(Y\le x\le Y^K).
\]
The unwanted twists can now be charged to \(\Sigma(x)\). What remains is to read the
fixed average \(G_{\mathcal K}\) of Definition~\ref{def:family-weights} from the local
prime sums bounded by \(E\).

\begin{construction}
\begin{tabular}{@{}l@{\hspace{1.2em}}l@{}}
\toprule
\textit{Object} & \textit{Operation or estimate} \\
\midrule
\(\sum_{j,t}\Delta_{j,t}(u)\) & Bound by \(E\) for every \(u\in A_0\) \\
Character \(\omega\) of \(H\) & Average over \(0\times A_0\) \\
\(\omega|_{0\times A_0}\ne1\) & Translate average vanishes \\
\(\omega|_{0\times A_0}=1\) & Keep \(\mathcal K\); bound the omitted sums by \(\Sigma(x)\) \\
\bottomrule
\end{tabular}
\caption{The character projection at a fixed scale. The local prime sums use the same integral weights at every occurrence of a given translate.}\label{fig:analytic-projection}
\end{construction}

\paragraph{Uniform estimates for the small-weight terms.}
\label{par:thin-host}
First isolate the terms that will disappear within the error allowance. A small weight
sends the argument of its prime sum far beyond \(x\). The family of characters at this step grows with \(x\).

\begin{lemma}[A uniform bound for character sums at large arguments]
Let \(N\) be the modulus defining the auxiliary field from
Section~\ref{sec:pair-selection}, so that
\(\log N\ll_\psi x^{\epsilon/5}\log^2 x\), and put
\(X'=\exp(x^{\epsilon/3})\). For each \(c>0\), uniformly over
all primitive characters \(\chi\) induced from this auxiliary
Galois group, one has
\begin{equation}\label{eq:thin-host}
    |G_\chi(y)|\le y\exp(-c x^{\epsilon/8})
                  \qquad (y\ge X'/3)
\end{equation}
once \(x\) in the selected window exceeds a bound depending only on
\(\psi,\epsilon,c\) and the fixed interval exponent.
\end{lemma}

\begin{proof}
\proofstep{Use a uniform zero-free region.}
The classical zero-free region for primitive Dirichlet characters of
conductor \(f\) excludes zeros with
\(\RePart s\ge1-c_0/\log(3f(2+|\Im s|))\), with the possible
exception of one real zero of a nonprincipal real character.
This region, the bound for the number of zeros in an interval of unit height, and the partial fraction
expansion of the logarithmic derivative are used with absolute
constants. For a real character the dyadic component factors through the units modulo eight. Indeed each unit congruent to one modulo eight is a square in \(\mathbb Z_2^\times\), and for a real character \(\chi_2\),
\[
 \chi_2(u^2)=\chi_2(u)^2=1\qquad(u\in\mathbb Z_2^\times).
\]
\cite{classical-region} Away from \(q,d\) the
remaining odd conductor divides the fixed odd part of the target
modulus. The two optional components at \(q,d\) are Legendre
characters by cyclicity. 

If a real character contains the quadratic factor at \(q\) or \(d\),
the zero-free rectangles of Section~\ref{sec:pair-selection} give the following
bound for a possible exceptional real zero \(\beta\):
\[
 L(\beta,\chi)=0\quad\Longrightarrow\quad \beta\le 1-1/400 .
\] The other real characters
belong to a fixed finite list. Nonvanishing at one for this list,
with the pole removed for the zeta function, gives a fixed zero-free
neighborhood of one. Decreasing \(c_0>0\) accordingly, the
stated region holds throughout this family without an exception.

\proofstep{Move the smoothed Perron integral.}
Put \(v=\log y\), \(T=\exp(\sqrt v)\) and
\(L=\log(3N(T+4))\). In the rectangle
\(1-c_0/(2L)\le\sigma\le1+1/v\), \(|t|\le T\), the logarithmic
derivative after removal of the pole of the zeta function is
\(O(L^2)\). For the zeros \(\rho\) of the character under consideration in the band \(|\Im\rho-t|\le1\), the zero-free region gives \(|\sigma+it-\rho|\gg L^{-1}\). The unit-band zero count now gives
\[
 \sum_{|\Im\rho-t|\le1}\frac{1}{|\sigma+it-\rho|}
 \ll L\,\#\{\rho:|\Im\rho-t|\le1\}\ll L^2 .
\] Comparing
the partial fractions at \(\sigma+it\) and \(2+it\) sums the other
bands with inverse-square weights, at cost \(O(L)\).


The smoothed Perron integral of
\[
       -\frac{L'}{L}(s,\chi)
                      \frac{y^{s-1}}{s(s-1)}
\]
on \(\RePart s=1+1/v\), with measure \(ds/(2\pi i)\),
equals \(S_\chi(y)/y\). Termwise inversion gives
\((1/n-1/y)_+\), including value zero at \(n=y\).
Absolute convergence and \(-\zeta'/\zeta(1+1/v)=O(v)\) bound
the tails above \(T\) by \(O(v/T)\). Put \(\sigma_*=1-c_0/(2L)\) for the new vertical line.
The residue at \(1\) is \(-L'(1,\chi)/L(1,\chi)\)
when \(\chi\) is nonprincipal. For the zeta function,
in the coordinate \(u=s-1\) the contribution is the coefficient of
\(u^{-1}\) in
\[
     \left(u^{-1}-\gamma+O(u)\right)
             \frac{e^{v u}}{u(1+u)}.
\]
It is \(v-1-\gamma\). Subtracting the respective residue
gives \(G_\chi(y)/y\). 

At height \(T\) each horizontal
segment has bounded length and the factor \(1/(s(s-1))\)
is \(O(T^{-2})\), so the horizontal error is \(O(L^2/T^2)\).
On the new line the pole term of \(-\zeta'/\zeta\) is \(O(L)\).
For \(|t|\le1\) one has
\[
 \frac{1}{|\sigma_*+it|\,|\sigma_*-1+it|}
        \ll \frac{1}{\sqrt{(c_0/(2L))^2+t^2}},
\]
whose integral is \(O(1+\log(2L))\).
For \(|t|\ge1\) the integrand is \(O(t^{-2})\).
The exponential factor on this line is \(e^{-c_0 v/(2L)}\).
With the log-derivative bound and the two original tails,
this gives
\[
 \frac{|G_\chi(y)|}{y}
  \ll v/T + L^2/T^2 + L^3 e^{-c_0 v/(2L)} .
\]

\proofstep{Compare the conductor with the argument.}
Since \(v\ge x^{\epsilon/3}-\log 3\), the conductor bound gives
\(\log(3N)\ll v^{3/5}\log^2(2v)\). Thus
\(L\ll v^{49/80}\) uniformly, while \(31/80>3/8\).
The exponential decay in \(v^{31/80}\) dominates every fixed
multiple of \(x^{\epsilon/8}\) in the exponent, including the
polynomial prefactors. This proves \eqref{eq:thin-host} using the
same primes as the cohomology window.
\end{proof}

\paragraph{The centered sum on the finite group.}
With the small weights under control we return to the whole finite group. Its
characters will separate the contributions that survive the translation average.

\begin{established}[Arguments of the inertia average]\label{hyp:inertia-fourier}
Let \(P\subset H\) be the index-two subgroup of Definition~\ref{def:translate-scale},
\(f_P=|P|=2n\), and \(\omega_P\) its quadratic character. For any lift
\(j\), top \(t\in T\) and translate \(u\in A_0\), translation pairs
\(H\setminus P\) with \(P\). On the selected chain put \(a=a_t\)
and let \(w_g\) be the weight at the image of \(g\). Uniformly for large \(x\),
\[
                0<w_g\le a,\qquad a/w_g\ge1\quad(g\notin P).
\]
Write \(\chi_\omega\) for the primitive inducer of the inflation of
\(\omega\in\widehat H\) by the Artin map.
\end{established}
\begin{proof}
Translation by an element outside \(P\) interchanges the two cosets.
The two ranges in Lemma~\ref{hyp:translate-weights} are positive.
Since \(\log\eta\le x^{\epsilon/4}\), \(b=e^{x^\epsilon}\) and
\(x/X'\to0\), both upper endpoints are at most \(2b\) eventually.
The top interval has \(a_t\ge(x-1)b\), so the same threshold serves all indices.
\end{proof}

\begin{lemma}[Fourier expression with inertia]
With the finite group and weights of Lemma~\ref{hyp:inertia-fourier},
for the stated lift \(j\), top index \(t\) and translate \(u\), one has
\begin{equation}\label{eq:color-balance}
  \Delta_{j,t}(u)=
  \RePart\frac1{f_P}\sum_{g\notin P}\sum_{\omega\in\widehat H}
       \overline{\omega(g)}\,w_g G_{\chi_\omega}(a/w_g).
\end{equation}
\end{lemma}

\begin{proof}
\proofstep{Invert the inertia averages.}
For a prime \(p\), let \(I_p\) be its
inertia subgroup and let \(F_p\) be a lift of its Frobenius in
\(H\). For \(\nu\ge1\) one has
\[
 \chi_\omega(p^\nu)
   =\omega(\nu F_p)\frac1{|I_p|}\sum_{v\in I_p}\omega(v).
\]
If \(\omega|_{I_p}\ne1\), the mean is zero, as is the value
of the primitive character at the ramified prime. If
\(\omega|_{I_p}=1\), the first factor is well defined on the
Frobenius coset and equals \(\chi_\omega(p^\nu)\). Thus
\[
 \frac1{f_P}\sum_{\omega\in\widehat H}
       \overline{\omega(g)}\chi_\omega(p^\nu)
 =\frac{|H|}{f_P|I_p|}\,\mathbf 1_{g\in\nu F_p+I_p}
 =\frac{2}{|I_p|}\,\mathbf 1_{g\in\nu F_p+I_p}.
\] Multiplied by
\(\log p\,(a/p^\nu-w_g)_+\) and summed over the prime powers
\(p^\nu\) and elements \(g\notin P\),
this gives the summand at lift \(j\) and index \(t\) in the arithmetic count.


At primes split in \(\mathbb Q(\sqrt{-d})\) the entire coset
lies in \(P\). At \(d\), write \(\nu F_d=(i,j,s)\). Definition~\ref{def:translate-scale} gives
\[
 (\nu F_d+I_d)\setminus P
   =\{(i,j,s),(i+1,j,s)\}\setminus P=\{(1,j,s)\}.
\]
The last two coordinates are independent of the choice of the lift \(F_d\);
the multiplicity is \(2/|I_d|=1\). At the other
unramified primes each contributing Frobenius element has weight two.
One may use any cutoff \(N_*\ge\max_g a/w_g\), since all terms
beyond it vanish. Interchange the finite sums. On the character side, factor the positive weight \(w_g\) out of the coefficient of \(\Lambda(n)\chi_\omega(n)\):
\[
 (a/n-w_g)_+
 =w_g\left(\frac{a/w_g}{n}-1\right)_+ .
\]
The inner sum is therefore the uncentered sum at \(a/w_g\).

\proofstep{Compute the two centering terms.}
The principal character contributes to the centering by
\[
 a\left(\log a-1-\gamma-f_P^{-1}\sum_{g\notin P}\log w_g\right).
\]
For a nonprincipal character the difference between the
uncentered and centered terms before weighting is
\(-aL'(1,\chi_\omega)/L(1,\chi_\omega)\) at each \(g\notin P\).
The character coefficient sums to zero except for
\(\omega_P\), where it sums to \(-f_P\).
Write \(\ell_P\) for the real logarithmic derivative
\(\RePart(L'(1,\chi_{\omega_P})/L(1,\chi_{\omega_P}))\).
The contribution to the centering term is \(+a\ell_P\).
Subtracting both contributions gives \eqref{eq:color-balance}.
\end{proof}

\paragraph{Replacing the weights before projection.}
The projection is now explicit. To compare it with the prescribed character average at
\(x\), we still have to account for the actual arguments at which the individual
weights are evaluated.

\begin{established}[Ranges of the two translated weights]\label{hyp:translate-weights}
Use the selected auxiliary field of \eqref{eq:thin-host} and the scale
of Lemma~\ref{lem:rescale-host-budget}. Put \(X'=e^{x^{\epsilon/3}}\),
\(b=e^{x^\epsilon}\), \(b'=xb/X'\), \(e_1(s)=((1,0),s)\), and
\(B_{t,u}=\{e_1(s+u-t):s\in T_2\}\). For every selected lift, top and
translate,
\[
 |a_t/b-x|\le1,\qquad
 \begin{cases} b\le w_g\le b+\eta+1,&g\in B_{t,u},\\
                0<w_g\le2b',&g\in(H\setminus P)\setminus B_{t,u},
 \end{cases}
 \qquad\log\eta\le x^{\epsilon/4}
\]
at sufficiently large scale. Here \(x^{2+8\epsilon}\le n\le2x^{2+8\epsilon}\)
and \(f_P=2n\).
\end{established}
\begin{proof}
The top is chosen in the first interval of Section~\ref{subsec:frobenius-selection},
so \(0\le a_t/b-x\le x/(kx^{13})\).
For the signed rows of Construction~\ref{hyp:family-prescribed-realization},
write \(b_v\) for the second top weight before the stage offset. The three
integer windows in Lemma~\ref{lem:marked-signed-jumps} have centers
and upper ends as follows, where \(w''=b/(nX')\):
\[
 \begin{array}{c|c}
 \text{location before translation}&\text{range containing the weight}\\ \hline
 s\in T_2 &[b,b+\eta+1]\\
 \text{designated top} &[b_v+\lfloor b'/x\rfloor,b_v+\lfloor b'/x\rfloor+\eta+1]\\
 \text{other small position} &[w'',w''+\eta+1].
 \end{array}
\]

The floor permits an endpoint loss of one. Translation sends precisely
the first set of positions to \(B_{t,u}\).
The offsets have width \(\eta=2m\lceil x^{c_H}\rceil\), with \(c_H\)
fixed before the final scale. Thus \(\log\eta=O(\log x)\).
The interval for \(b_v\) in Section~\ref{subsec:frobenius-selection}
has lower end \(b'\) and relative width \(2n\delta x^\epsilon\Lambda/k\),
which tends to zero in the selected range. Also
\[
 (1+\eta)/b'\to0,\qquad w''/b'=1/(nx),\qquad
                  \log\eta\le x^{\epsilon/4}.
\]
Both complementary windows consequently have upper end at most \(2b'\)
for all sufficiently large \(x\).
The period and level identifications in Proposition~\ref{prop:translate-completion}
preserve these integral weights.
\end{proof}

\begin{proposition}[Averaging over translates]\label{prop:translate-fourier}
Use the character-weight data of Lemma~\ref{hyp:inertia-fourier}
with the translate weights of Lemma~\ref{hyp:translate-weights}
at \(x\) in a selected window. Suppose that
\(\sum_{j,t}\Delta_{j,t}(u)\le E\) for every \(u\in A_0\).
For \(x\) above one threshold independent of \(u,t,j\) and the
auxiliary field in that window, one has
\begin{equation}\label{eq:fourier-compare}
 \frac{Q\delta_2}{2}
       \bigl(G_{\mathcal K}(x)-\Sigma(x)-4x^{1/2+\epsilon}\bigr)
       \leq E+\frac{2Q}{x}.
\end{equation}
\end{proposition}

\begin{proof}
\proofstep{Replace the large weights and discard the small weights.}
For the group elements with larger weights, change the weight from \(w_g\) to \(b\).
The positive-part function is \(1\)-Lipschitz and both weights
are at least \(b\). Consequently,
\[
 \begin{aligned}
 |w_g S_{\chi_\omega}(a/w_g)-b S_{\chi_\omega}(a/b)|
 &\le (w_g-b)\sum_{n\le a/b}\Lambda(n)
 \ll (x+1)(\eta+1).
 \end{aligned}
\]
The nonprincipal centering term is independent of \(w_g\).
For the principal term, the difference is bounded by
\(a|\log(w_g/b)|\le(a/b)(w_g-b)\). For a group element with a smaller weight,
\(a/w_g\ge (x-1)X'/(2x)\ge X'/3\) for \(x\ge3\), so its entire
contribution before multiplying by the character coefficient is at
most \((x+1)b\exp(-c x^{\epsilon/8})\) by \eqref{eq:thin-host}.


There are at most \(f_P\) terms in the \(g\)-sum and exactly
\(2f_P\) in the character sum. After dividing by \(f_P\), the
respective errors are bounded by
\[
 O(f_P(x+1)(\eta+1)),\qquad
 2f_P(x+1)b e^{-c x^{\epsilon/8}}.
\]
Divide these bounds by \(b/x\). For a fixed \(c>0\) they are, respectively,
\[
 \begin{aligned}
 &O\!\left(x^{4+8\epsilon}\exp(x^{\epsilon/4}-x^\epsilon)\right)=o(1),\\
 &O\!\left(x^{4+8\epsilon}\exp(-c x^{\epsilon/8})\right)=o(1).
 \end{aligned}
\]
Only \(x,\epsilon,c\) and the fixed character enter these upper bounds. Consequently the contribution from the
larger weights, denoted by \(\Theta_{u,t}\), satisfies
\(\RePart\Theta_{u,t}\le\Delta_{j,t}(u)+2b/x\) for each lift.
Its formula is
\[
  \Theta_{u,t}=\frac b{f_P}\sum_{s\in T_2}\sum_{\omega\in\widehat H}
       \omega(-e_1(s+u-t))G_{\chi_\omega}(a_t/b).
\]

\proofstep{Project to the quotient characters.}
Use the group \(\widehat H_0\) of Definition~\ref{def:translate-scale}, consisting of the characters trivial on
\(0\times A_0\). For each \(s\in T_2\subset A_0\),
\[
 \frac1{|A_0|}\sum_{u\in A_0}\omega(-e_1(s+u-t))
   =\omega(e_1(t))\mathbf 1_{\omega\in\widehat H_0}.
\]
The average of \(\Theta_{u,t}\) is therefore
\[
 \frac{b\delta_2}{2}
 \sum_{\omega\in\widehat H_0}
      \omega(e_1(t))G_{\chi_\omega}(a_t/b),
 \qquad \frac{bn_2}{f_P}=\frac{b\delta_2}{2}.
\] 

The quotient by \(0\times A_0\) has order
\(4o\). Its characters come from powers of the fixed character
and their twists by the two quadratic characters of the prime
pair. Indeed imposing triviality on \(A_0\) removes the growing
real two-power factor. Write \(\eta_q,\eta_d\) for the quadratic characters of the auxiliary primes. The Chinese remainder theorem makes the target value and the two prime signs independent. On the quotient the characters thus have the parametrization
\[
 \psi^a\eta_q^b\eta_d^c,
 \qquad 0\le a<o,\qquad b,c\in\{0,1\},
\]
where the product denotes the associated Artin character. For the powers of the target,
the restriction to \(A\) is the associated character of \(A/A_0\):
one can prescribe the target value and adjust both prime signs
independently. Also \(e_1(0)\) generates inertia at \(d\), where
the target is unramified. Thus its phase is one even for an odd
target. 

For the remaining characters trivial on \(0\times A_0\), let \(c_\chi\) denote the phase coefficient. Its modulus is at most one, so taking real parts gives
\[
 \sum_{\chi\ {\rm remaining}}\Re\!\left(c_\chi G_\chi(x)\right)\ge-\sum_{\chi\ {\rm remaining}}|G_\chi(x)|=-\Sigma(x).
\]

\proofstep{Average over the cosets represented in \(T\).}
The proportions of \(T\) in the identity and other cosets of
\(A/A_0\) are \(4/(3o+1)\) and \(3/(3o+1)\), respectively,
each up to \(x^{-1+\epsilon}\) by Section~\ref{subsec:frobenius-selection}.
For a nontrivial character \(\xi\) of \(A/A_0\), weight each of these
proportions by the value of \(\xi\) on the corresponding coset. This gives
\[
 \frac1{n_T}\sum_{t\in T}\xi(t)
  =\frac4{3o+1}+\frac3{3o+1}\sum_{c\ne A_0}\xi(c)
     +O(o x^{-1+\epsilon})
  =\frac1{3o+1}+O(o x^{-1+\epsilon}).
\]
The second equality uses \(\sum_c\xi(c)=0\). The mean for the
trivial character is exactly one.

For the fixed nonprincipal target characters the centered sums
are \(O_\psi(x\log^2(2x))\) by their defining series.
For all characters trivial on \(0\times A_0\), replacing \(a_t/b\) by \(x\) introduces an error of
\(O_\psi(\log^2(2x))\), because for \(1\le y,y'\le x+1\) the positive-part kernel gives
\[
 |S_\chi(y)-S_\chi(y')|
 \le |y-y'|\sum_{n\le x+1}\frac{\Lambda(n)}n
 \ll |y-y'|\log^2(2x).
\]
The logarithmic derivatives are fixed on the target family and
bounded by \(4\log x\) for the remaining characters in this sum. The elementary
principal correction has derivative \(O(\log x)\). The total of
these errors after averaging over the top indices is
\(O_\psi((1+x^\epsilon)\log^2(2x))\), at most
\(4x^{1/2+\epsilon}\) above a fixed threshold.

If \(\mathcal A=m\sum_{t\in T}\RePart\bigl(|A_0|^{-1}
                  \sum_{u\in A_0}\Theta_{u,t}\bigr)\),
the character calculation bounds \(\mathcal A\) below by the left
side of \eqref{eq:fourier-compare}. 

On the other hand,
\[
 \mathcal A\le \frac1{|A_0|}\sum_{u\in A_0}
                \sum_{j,t}\Delta_{j,t}(u) + \frac{2mn_Tb}{x}
             \le E+\frac{2Q}{x}.
\]
Both bounds refer to the same translate family.
\end{proof}

Projection has identified the desired characters. Lemma~\ref{lem:comparison} first divides by the total weight \(Q=n_Tmb\), then by the surviving proportion \(\delta_2/2\). The resulting error tells us how strong the two lattice estimates on these chains must be.


\begin{lemma}[Normalization of the comparison]\label{lem:comparison}
Let \(x\ge2\) and \(\epsilon>0\). The parameters
\(n,m,n_T,n_2,b,k\) are positive. Let \(a_E>0\) be a size parameter
and let \(C\ge0\) be the exponent coefficient in the error bound.
Assume
\(\Sigma(x)\le4o\,x^{1/2+\epsilon}\). Suppose
\eqref{eq:fourier-compare} holds and
\[
 E\le x^{C\epsilon}
 \left[n_T a_E\left(1+\frac{m^2}{k}\right)
      +m\left(2nmb x^\epsilon
               +\frac{n_2 k b x^{2\epsilon}}{n^2}\right)\right],
\]
and that the sizes satisfy
\[
\begin{gathered}
 a_E\le 2xb,\qquad n/n_T\le 2x^{\epsilon/8},\qquad
 x^{-\epsilon}\le\delta_2\le1,\\
 x^{1/2-\epsilon}\le m\le x^{1/2+\epsilon/2},\qquad
 xm\le k\le n^2 x^{1/2-\epsilon}.
\end{gathered}
\]
Then for \(x\) sufficiently large in terms of \(o,\epsilon\),
\[
 G_{\mathcal K}(x)\le x^{1/2+(C+4)\epsilon}.
\]
\end{lemma}

\begin{proof}
\proofstep{Normalize by total weight.}
Divide each summand in the bracket by \(Q=n_Tmb\). The
individual estimates are
\[
 \begin{aligned}
 \frac{a_E}{mb}\left(1+\frac{m^2}{k}\right)
    &\le\frac{2x}{m}+\frac{2xm}{k}
      \le 2x^{1/2+\epsilon}+2,\\
 2\frac{n}{n_T}m x^\epsilon
    &\le 4m x^{9\epsilon/8}\le4x^{1/2+13\epsilon/8},\\
 \frac{n}{n_T}\,\delta_2\,\frac{k}{n^2}x^{2\epsilon}
    &\le 2x^{1/2+9\epsilon/8}.
 \end{aligned}
\]
The last line uses \(\delta_2\le1\) and the upper bound on \(k\).
Their sum is at most \(10x^{1/2+13\epsilon/8}\).

\proofstep{Divide by the density of the retained weights.}
Divide \eqref{eq:fourier-compare} by \(Q\delta_2/2\).
Using \(\delta_2\ge x^{-\epsilon}\), we get
\[
 G_{\mathcal K}(x)-\Sigma(x)
 \le 4x^{1/2+\epsilon}
     +20x^{1/2+(C+21/8)\epsilon}
     +4x^{-1+\epsilon}.
\]
Since \(x\ge2\) and \(C\ge0\), the sum is at most
\(26x^{1/2+(C+21/8)\epsilon}\). Once \(x^{3\epsilon/8}\ge26\),
it is at most \(x^{1/2+(C+3)\epsilon}\). Restoring \(\Sigma(x)\)
gives the result once \(x^\epsilon\ge1+4o\).
\end{proof}

The two terms \(x/m\) and \(m\) in the normalized bound explain the
choice of \(m\) near \(\sqrt{x}\). The threshold in the lemma
is independent of the auxiliary field and the dimensions of the lattices.
Section~\ref{subsec:fixed-weight-absorption} obtains \(C=5\) when its divisibility and slope bounds hold in common at the selected points. This gives the exponent \(1/2+9\epsilon\).
The arithmetic conclusion needed for the analytic argument is the following bound at every translate.

\begin{windowtheorem}[Uniform arithmetic estimate for a primitive character]\label{hyp:uniform-arithmetic-data}
Fix a primitive character \(\psi\) of order \(o\). The range of the interval assertion is
\[
 \exists\epsilon_0(\psi)\in(0,1/52)\quad
 \forall\epsilon\in(0,\epsilon_0(\psi))\quad
 \forall Z\ge1\quad\exists Y>Z .
\]
For these choices the following data exist. There is an outer interval \([U,U^J]\) containing \([Y,Y^{3/2}]\). Choose one auxiliary prime pair there as in Proposition~\ref{prop:quadratic-prime-pair}, with the prime sizes measured at \(U\). On the smaller interval the field and numerical data may vary:
\[
                  Y\le x\le Y^{3/2}\quad\longmapsto\quad E(x).
\]
At each such \(x\), choose the data of Lemmas~\ref{hyp:inertia-fourier} and~\ref{hyp:translate-weights} at that pair, satisfying the size inequalities of Lemma~\ref{lem:comparison} and its bound for \(\Sigma(x)\).

With the notation of Definition~\ref{def:translate-errors} put
\[
 \mathcal B(x)=n_Ta_E(1+m^2/k)
       +m\left(2nmbx^\epsilon+n_2kb x^{2\epsilon}/n^2\right).
\]
On these data the centered local contributions obey the simultaneous bound
\begin{equation}\label{eq:uniform-raw-gap}
 \max_{u\in A_0}\sum_{j=1}^{m}\sum_{t\in T}\Delta_{j,t}(u)
       \le x^{5\epsilon}\mathcal B(x).
\end{equation}
All thresholds for the estimates used on the interval may depend on the fixed character and \(\epsilon\) and are uniform in its subsequent scale and translate choices.
\end{windowtheorem}

\begin{proof}
Fix the weights of Definition~\ref{def:family-weights} before choosing
\(\epsilon\). Take the minimum of \(1/53\) and the positive
small-parameter bounds in Lemma~\ref{lem:window-parameter-choice}.
Lemma~\ref{lem:disjoint-common-pair-windows} gives disjoint intervals
with left endpoints tending to infinity. At each of their scales
take the base data of Section~\ref{subsec:frobenius-selection}.
With \(\delta=x^{-\epsilon/8}\), Proposition~\ref{prop:top-set-cycles}
and Lemma~\ref{lem:second-set-progressions} give
\[
 n_T\ge\delta n/2,\qquad x^{-\epsilon}n\le n_2\le n .
\]

The coset proportions are those of \eqref{eq:coset-shares}.
The host Artin map induces precisely the characters of
Definition~\ref{def:translate-errors}. In the annihilator of
\(A_0\) there are at most \(4o\) characters. Each omitted one is
a primitive inducer of a power of \(\psi\) twisted by one of the
three nontrivial quadratic characters of the prime pair. Thus
Proposition~\ref{prop:quadratic-prime-pair} on the outer interval gives
\[
 \Sigma(x)\le4o\,x^{1/2+\epsilon/2}\le4o\,x^{1/2+\epsilon}.
\]

Take the common member for each translate from
Proposition~\ref{prop:all-window-applications}. The size choices
of Lemma~\ref{lem:rank-window} give the \(m\)-range after a fixed
increase of the lower threshold. With the top size \(a_E\le2xb\),
the other rank comparisons reduce to
\[
   xm\le k,\qquad k/n^2\le2x^{1/2-6\epsilon}
                           \le x^{1/2-\epsilon}.
\]
Lemma~\ref{lem:uniform-numerical-application} bounds every translate
on this family by the right side of \eqref{eq:uniform-raw-gap}.
Choose a left endpoint beyond \(Z\) and these common thresholds.
All the bounds then hold throughout the same closed interval.
\end{proof}

\begin{proof}[Deduction of Proposition~\ref{prop:arithmetic} from Theorem~\ref{hyp:uniform-arithmetic-data}]
For the chosen \(\psi,\epsilon\), increase the prescribed lower endpoint beyond the thresholds in Proposition~\ref{prop:translate-fourier} and Lemma~\ref{lem:comparison}. For each scale take \(E=x^{5\epsilon}\mathcal B(x)\) in the former. Inequality~\eqref{eq:uniform-raw-gap} verifies its hypothesis for every translate with this one value of \(E\). The lemma now applies at \(C=5\), giving
\[
   G_{\mathcal K}(x)\le x^{1/2+(5+4)\epsilon}
           \qquad(Y\le x\le Y^{3/2}).
\]
The lower endpoint could be chosen above any prescribed bound. This proves the quantified estimate \eqref{eq:arithmetic-window}.
\end{proof}

\begin{auxlemma}[A comparison at a maximizing translate]\label{lem:maximizing-translate}
Fix the finite group, character data, top weights, scale and coset
proportions of Proposition~\ref{prop:translate-fourier}. Keep these
data fixed as \(u\) varies. Before choosing any of the smaller integral
weights, form the real function on the finite group \(A_0\)
\[
 B_x(u)=\frac{mb}{f_P}\RePart\!
       \sum_{t\in T}\sum_{s\in T_2}\sum_{\omega\in\widehat H}
       \omega(-e_1(s+u-t))G_{\chi_\omega}(a_t/b).
\]
Choose any \(u_*\in\operatorname*{arg\,max}_{u\in A_0}B_x(u)\).
Suppose the weight inequalities of Lemma~\ref{hyp:translate-weights}
for the individual lifts, and the inertia data of Lemma~\ref{hyp:inertia-fourier},
are realized at \(u_*\). Under the threshold and conductor conditions
of Proposition~\ref{prop:translate-fourier}, it suffices for
\eqref{eq:fourier-compare} that
\[
             \sum_{j,t}\Delta_{j,t}(u_*)\le E.
\]
\end{auxlemma}
\begin{proof}
Only the base data enter \(B_x\). Averaging its explicit sum over \(A_0\)
kills exactly the characters nontrivial on \(0\times A_0\).
The coset-proportion calculation in Proposition~\ref{prop:translate-fourier}
therefore gives
\[
 \frac{Q\delta_2}{2}
   \bigl(G_{\mathcal K}(x)-\Sigma(x)-4x^{1/2+\epsilon}\bigr)
 \le |A_0|^{-1}\sum_u B_x(u) \le B_x(u_*).
\]
At \(u_*\), the large-weight replacement and the small-weight estimate
in that proposition use only the weights of one lift and top label at
a time. Each replacement costs at most \(2b/x\).
Summing over the \(mn_T\) pairs yields
\[
 B_x(u_*)\le\sum_{j,t}\Delta_{j,t}(u_*)+2Q/x\le E+2Q/x.
\]
All terms are on the same scale. The maximizer is chosen from the
Fourier sum; its location need not agree with a translate selected
independently for geometric reasons.
\end{proof}

\begin{auxlemma}[Recovering the interval estimate from square arguments]\label{lem:square-window-transfer}
Fix a primitive target \(\psi\). Write \(G=G_{\mathcal K(\psi),w_\psi}\)
as in Definition~\ref{def:family-weights}. Suppose that for each
sufficiently small \(\eta>0\) there are arbitrarily large \(Y\) such that
\[
       G(k^2)\le (k^2)^{1/2+9\eta}
       \qquad(k\in\mathbb N,\ Y\le k^2\le Y^{3/2}).
\]
Then for every sufficiently small \(\epsilon>0\) and arbitrarily large \(Y\),
\[
           G(y)\le y^{1/2+9\epsilon}\qquad(Y\le y\le Y^{3/2}).
\]
\end{auxlemma}
\begin{proof}
Use the square bound with \(\eta=\epsilon/2\).
Every \(y\) in such an interval has a square \(z\) in the same interval
with \(|z-y|\le 3\sqrt y\) once \(Y\) is large. Indeed, one of
\(\lfloor\sqrt y\rfloor^2\) and \(\lceil\sqrt y\rceil^2\) lies inside:
if both were outside, the interval would contain no square. Yet
\(\sqrt{Y^{3/2}}-\sqrt Y>1\) for all sufficiently large \(Y\).

For \(y,z\ge2\), put \(v=\max(y,z)\). The function
\(t\mapsto(t/n-1)_+\) is \(1/n\)-Lipschitz; it vanishes throughout
the segment between \(y\) and \(z\) for \(n>v\).
The factorial bound \eqref{eq:factorial-mangoldt-bound} and the derivative
of the principal centering term give
\[
 |G(y)-G(z)|\ll_\psi |y-z|\log(2v).
\]
The nonprincipal centering coefficients are constants on this fixed
finite family. Since \(z\le 2y\) for large \(Y\), we have uniformly
\[
 G(y)\le 2^{1/2+9\eta}y^{1/2+9\eta}
              +O_\psi(\sqrt y\log(2y))\le y^{1/2+9\epsilon}.
\]
The last inequality follows after division by \(y^{1/2+9\epsilon}\),
since \(\epsilon-\eta>0\).
\end{proof}

In particular one may fix the prime pair for a late interval at error
parameter \(\eta\), and work only at the finitely many squares in it.
At each square, fix the base data of Lemma~\ref{lem:maximizing-translate}
first. It suffices to realize both divisions and their uniform numerical
and Gram bounds on one completed member at the distinguished translate:
\[
 x=k^2\ \longmapsto\ B_x\ \longmapsto\ u_*(x)
   \ \longmapsto\ P(x,u_*) .
\]
The calculation in the proof of Corollary~\ref{cor:raw-gap-completed-scale}
on this member and the normalization in Lemma~\ref{lem:comparison} give
\[
 \sum_{j,t}\Delta_{j,t}(u_*)\le x^{5\eta}\mathcal B(x)
       \quad\Longrightarrow\quad G_{\mathcal K}(k^2)\le(k^2)^{1/2+9\eta}.
\]
Lemma~\ref{lem:square-window-transfer} gives
\eqref{eq:arithmetic-window}. A single list of allowed word primes at
each such translate suffices if it supports all these comparisons.
Showing that such a member exists with the required constants is still
an arithmetic problem.

\paragraph{What remains to construct on an interval.}\label{par:arithmetic-construction-inputs}
The numerical argument is complete once the local prime sums satisfy
\eqref{eq:uniform-raw-gap} or its square-argument variant. To reach it by
theta lifting we have to interpret these sums as degrees of lattices.
We follow the construction at one translate.

At a fixed scale and translate, assign the split word primes. The two CM types induced from the quadratic subfield have principal rank-one sources by Lemma~\ref{lem:marked-principal-cm-sources}. Tensoring with the Hermitian lattice gives the definite points of Lemma~\ref{lem:marked-definite-cm-points}. At a positive-dimensional vertex the rational cusp sections also define the integral ball. If \(s=\sum_i z_i b_i\) in a basis of sections, then
\[
    z\in B_v\quad\Longleftrightarrow\quad
    \left|\sum_i z_i b_i(y)\right|_v\le1
           \quad\text{for every admissible integral test }y .
\]
This is the same norm in both presentations of the row (Lemma~\ref{lem:constructed-signed-coefficient-row}). A finite set of CM evaluations gives a weaker but quantitative assertion: if \(A\) is their invertible matrix and \(\pi^c A^{-1}\) is integral, then \(Ax\in\pi^{-b}R^d\) gives \(x\in\pi^{-(b+c)}R^d\). At good ordinary level Lemma~\ref{lem:ordinary-hecke-evaluations} gives detection modulo \(p\) on an entire prime-to-\(p\) Hecke orbit.

The initial signed tests choose the characters. Only the reference system is copied coefficientwise. On a new tuple the specialized columns span the good row, so their adjoint detects each eigenspace for the transferred operators (Lemma~\ref{lem:new-ray-type-occurrence}). Thus Construction~\ref{hyp:family-prescribed-realization} gives the order of choice
\[
 (K_s,K_0,\text{reference row})\ \longmapsto\ (\varphi_i)_i
     \ \longmapsto\ P'\ \longmapsto\
        (\text{cover at }p_0,\text{cover at }p_1).
\]
For a copied character, \(\lambda_w(T)=j(\lambda_v(T))\) at every tested unramified Hecke operator. A newly occurring system satisfies its transferred equalities on the specified algebra. The two covers in the last step both belong to \(P'\).

The height estimate starts with a period on that row.
Lemma~\ref{lem:weighted-fitted-translate} provides a bounded translate
with nonzero projection of controlled size. The resulting level and
the lattice estimate \eqref{eq:fitted-auxiliary-depth} make the orbital
tail strictly smaller than the stabilizer term. The core span and the
successive Fourier restrictions in Lemmas~\ref{lem:fitted-core-span}
and~\ref{lem:fitted-fixed-ascent} keep the same bottom row while choosing
the tests above it. Consequently Lemma~\ref{lem:fitted-completed-nonvanishing}
proves pointwise nonvanishing at the prescribed depth.

The cohomological division uses the integral cyclotomic module
of Proposition~\ref{prop:ramified-integral-wach-comparison} and the
restricted cocycles of Lemma~\ref{lem:ramified-restricted-herr} over the
same ramified field. For a crystalline lattice with jumps in \([-a,b]\)
the two height bounds read
\[
             Q^a\varphi^*N\subset N,\qquad Q^b N\subset\varphi^*N .
\]
The scalar period at \(2\) can be computed on the cyclotomic tower. For
\(F_0/\mathbb Q_2\) unramified and \(F_s=F_0(\mu_{2^s})\),
Lemma~\ref{lem:dyadic-integral-period} constructs integral cohomology
classes whose dual exponentials give
\[
 2^{s-1-e-v_2((e-1)!)}\mathcal O_{F_s}
     \subseteq 2^{1-e-v_2((e-1)!)}\mathcal O_{F_s}
     \subseteq\exp^* H^1(F_s,\mathbb Z_2(-e))
          \qquad(e>0,\ s\ge2).
\]
The restriction from \(E\) to \(E(i)\) in Lemma~\ref{lem:dyadic-host}
is a different calculation. Its possible loss of one unit of length per twist
is proved by inflation and restriction.

On the geometric side only the degree-two part of the relative
Frobenius comparison enters the mixed intersection:
\[
             N_2(FX-X)=N_2(D_p).
\]
On a framed chart Lemmas~\ref{lem:flag-local-cyclotomic-chain} and~\ref{lem:flag-local-hodge-quotients} identify the two Frobenius quotients and their Hodge end maps. Lemma~\ref{lem:flag-finite-frobenius-model} descends this diagram to one finite model, including its boundary. At a prime of the different, the integral eigenlattices need not split the original Tate lattice. Proposition~\ref{prop:flag-cutoff-test} compares their flag metrics by successive changes of lattice. The changes of the graded squares cancel in the full sum; polarization then identifies the contribution from the selected half. This gives the input at ramified primes for the slope calculation in Section~\ref{sec:flag-wall-slope}.

At the definite endpoint,
Lemma~\ref{lem:definite-schur-gram} bounds the Gram determinant and
evaluation of the standard Schur lattice. The definite component has no boundary. Lemma~\ref{lem:terminal-valuative-realization} identifies the integral fibre on the geometric row and gives zero for the degree \(d_Z\) of the discrepancy in \eqref{eq:terminal-theta-frame-defect}. With \(g=fr\) and \(I_Z\) the index of the CM eigenlattices, this gives
\[
 \sum_\tau c_{\tau,Z}-2t_Z= I_Z-\tfrac g2\log(2\pi),\qquad
                   n\Delta_Z\le2U .
\]
The separate displacement of the weights is bounded in Lemma~\ref{lem:terminal-centred-weights}.

Finally the real integral must use the phases of the same finite and compact types. Lemma~\ref{lem:oriented-terminal-fibre} shows how the terminal frame determines the matrix test after integration along the two stabilizers. The cell calculation and the full Gram polynomial then give integrable bounds in Lemma~\ref{lem:empty-word-exhaustion}. Lemma~\ref{lem:residual-double-cofactor} identifies the double cofactor with the Gram determinant times a compact Capelli eigenvalue. For the two singly contracted cofactors,
Lemma~\ref{hyp:selected-orbital-coefficient-bounds} computes the
operators in the actual compact lines. Flattening the radial Haar
measure and changing from the positive boost to the oriented frame
give the phases in \eqref{eq:three-residual-cofactors}.

At a completed scale Corollary~\ref{cor:raw-gap-completed-scale}
combines the two index estimates. Passing from that scale to
Theorem~\ref{hyp:uniform-arithmetic-data} uses the following order:
\[
 \text{fix the family on }S\ ;\quad
 \text{take its common threshold }X_*\ ;\quad
 Y_\nu>\max\{Z,X_*\}.
\]
The field may change with the dyadic index inside the chosen
interval. Its prime pair, the weights on the target characters,
and the comparison constants continue to serve all its scales.
Lemma~\ref{lem:disjoint-common-pair-windows} proves the assertion
about the pair. Lemma~\ref{lem:uniform-numerical-application}
computes the last threshold from the estimates on this family.
